\section{Introduction}\label{sec:introduction}

An electronic ground-energy problem specifies both an external potential
and the states over which the energy is minimized. These choices matter
for computational hardness. A finite collection of orbitals can encode a
hard spin system, yet an arbitrary continuum state may have lower energy
than every state in that collection. Likewise, a freely designed
potential can create a useful array of wells without being the potential
of positive point nuclei. We prove hardness when the only external data
are rational positions of nuclei of charge one, and the minimization is
over the full spinful fermionic continuum space.

Schuch and Verstraete proved QMA-completeness for a Hubbard model with
local magnetic fields and developed a continuum realization using
designed scalar and spin-dependent
potentials~\cite{SchuchVerstraete2009}. Their reduction to Heisenberg
exchange and their comparison with the complementary one-particle space
are predecessors of the finite and continuum estimates used here.
Requiring the external potential to come solely from positive point
nuclei imposes a further restriction.

O'Gorman, Irani, Whitfield, and Fefferman proved QMA-completeness for
electronic structure in a supplied finite orbital
basis~\cite[Theorem 1]{OGorman2022}. Their journal version includes a
construction with positive unit point
charges~\cite[Appendix D, Corollary 1]{OGorman2022}. The supplied basis
restricts the admissible many-electron states, as the authors emphasize
in their discussion of the infinite-dimensional problem in Section~V.
A projected energy supplies a variational upper bound on the continuum
infimum; transferring a NO instance also requires a lower bound on every
omitted continuum state. The present theorem retains both unit nuclei
and unrestricted spinful continuum minimization.

\subsection{The electronic promise problem}

Let $M\ge1$ clamped nuclei have pairwise distinct positions
$R_1,\ldots,R_M\in\mathbb Q^3$, and let $N\ge1$ be the number of
electrons. Each nucleus has charge $Z_\alpha=1$. In atomic units the
electronic Hamiltonian is
\begin{equation}\tag{H}\label{eq:electronic-H}
 H=-\frac12\sum_{i=1}^N\Delta_{x_i}
   -\sum_{i=1}^N\sum_{\alpha=1}^M\frac1{|x_i-R_\alpha|}
   +\sum_{1\leq i<j\leq N}\frac1{|x_i-x_j|}.
\end{equation}
Its Hilbert space and form domain are
\[
 \mathcal H_N=\bigwedge^N\bigl(L^2(\mathbb R^3)\otimes\mathbb C^2\bigr),
 \qquad
 \mathcal Q_N=H^1(\mathbb R^{3N};(\mathbb C^2)^{\otimes N})
                  \cap\mathcal H_N.
\]
Antisymmetry exchanges position and spin together. No spin sector is
selected. We use the lower-bounded self-adjoint operator associated with
the closed Coulomb quadratic form on $\mathcal Q_N$, and write
$E_0=\inf\operatorname{spec}H$. Equivalently, $E_0$ is the infimum of
this form over normalized vectors in $\mathcal Q_N$, whether or not it
is attained; see~\cite[Theorems 2.13 and 2.19]{Teschl2009} for the
closed-form representation and variational characterization.
Section~\ref{sec:many-body} gives the form bounds and core argument
used in the continuum comparison.

\begin{table}[htbp]
\centering
\small
\begin{tabular}{@{}p{.23\linewidth}p{.70\linewidth}@{}}
\toprule
Input or convention & Specification \\
\midrule
Nuclear positions & Pairwise distinct rational triples $R_\alpha$ \\
Nuclear charges & $Z_\alpha=1$ for every $\alpha$ \\
Electron number & $N\ge1$, encoded in unary \\
Thresholds & Rational numbers $a<b$ \\
Rational encoding & Signed binary numerator and positive binary denominator
in lowest terms; fixed self-delimiting encoding \\
Input length & $L$: total length of the binary and unary data \\
Promise & $b-a\ge L^{-1}$; YES if $E_0\le a$, NO if $E_0\ge b$ \\
Energy convention & Electronic energy; clamped nuclear repulsion omitted \\
\bottomrule
\end{tabular}
\caption{The unit-charge input model. The decision-threshold separation is
an input precision promise, not a spectral-gap assumption.}
\label{tab:input-model}
\end{table}

The instance and its encoding are specified in Table~\ref{tab:input-model}.
The promise excludes energies strictly between the thresholds. Nuclear
repulsion is omitted: for clamped positions it is a fixed
geometry-dependent scalar, whereas~\eqref{eq:electronic-H} defines the
electronic-energy convention used throughout.

\begin{theorem}[Unit-charge Coulomb hardness]\label{thm:main}
The electronic promise problem for~\eqref{eq:electronic-H} on
$\mathcal H_N$, with distinct rational nuclear positions, all nuclear
charges equal to one, unary electron number, and rational thresholds
separated by at least $L^{-1}$, is QMA-hard under deterministic classical
polynomial-time many-one reductions. The reduction produces
polynomially many nuclei and electrons, polynomial rational bit lengths,
and threshold separation at least one. The only output data are the
nuclear positions, electron number, and thresholds.
\end{theorem}

Thus the bounded-charge problem holds with the absolute charge bound
$Z_{\max}=1$ and precision exponent $c=1$. The theorem imposes no
neutrality, binding, existence of a ground eigenvector, molecular
spectral-gap, or lower nuclear-separation promise. It supplies no orbital
basis, electronic state, magnetic field, or additional external
potential. Membership in QMA is not asserted. The constructed inequality
$b-a\ge1$ is stronger than the input requirement $b-a\ge L^{-1}$; neither
inequality specifies an excitation gap of the molecule.

The strict-binding excess-charge bound
in~\cite[Theorem 1.1]{OpenAIIonization2026} bounds how many electrons
fixed nuclei can bind strictly. Such a bound supplies neither
a finite-dimensional truncation nor the continuum lower bound required
here.

The companion article~\cite[Theorem 1.1 and Section 2]{OpenAIBinary2026} develops the
finite-gadget framework in greater detail and proves full-continuum
hardness with binary-encoded nuclear charges. Its large-charge analytic
construction and the polynomial-scale realization here are distinct.
Our proof starts directly from the same external spin source and
includes the finite estimates it uses; it does not invoke the companion's
hardness theorem.

\subsection{Proof strategy}

The starting point is the QMA-hard energy problem for signed Heisenberg
exchange on finite square-lattice subgraphs. We use the spatially sparse
construction of Cubitt, Montanaro, and Piddock, whose polynomial
interaction scales are essential here~\cite[Theorem 48, Lemma 47, and
Section 10.1]{CMP2019}. Section~\ref{sec:source} puts this source in a
rational encoding with an explicit remaining promise gap.

The finite part of the reduction has two stages.
Section~\ref{sec:spin} replaces signed exchange by positive exchange
using singlet-pair mediators~\cite[Section 2.4 and Lemma 5]{PiddockMontanaro2017},
and gives a planar layout with independently adjustable link lengths.
Section~\ref{sec:wells} associates one localized
spatial mode, with two spin states, to each of the resulting $m$ sites.
With $m$ electrons, a positive occupancy penalty favors one electron per
site. Virtual hopping through double occupancy then gives the prescribed
positive spin exchange, by the superexchange mechanism of
Anderson~\cite{Anderson1959}. The direct Coulomb coefficients enter that
penalty and its virtual-excitation denominators explicitly.

The remaining construction realizes these modes using unit nuclei. We
first work in dilated coordinates $y=\lambda x$, with $\lambda$ a
polynomially large scale. After dividing energies by $\lambda^2$, a unit
nucleus has coupling $1/\lambda$. Section~\ref{sec:density} constructs a
nonnegative continuous nuclear density producing the planar wells and
transverse confinement. A broad, constant-density slab supplies the
positive curvature needed to keep the modified density nonnegative.
Small changes in the well depths supply a one-site counterterm for the
long-range direct electron repulsion.

Section~\ref{sec:nuclei} replaces the density by distinct point nuclei.
A smooth flow transports uniform density to the constructed density.
Applying this flow to a tensor-product two-point Gauss rule preserves
equal node masses. The mass is chosen to be exactly $1/\lambda$, so
returning to physical coordinates gives charge one at every nucleus.
Rational rounding of the positions is included in the estimates.

The key analytical distinction is between the replacement error on
arbitrary form-domain vectors and its compression to the localized
modes. The latter is smaller because the modes have smooth Coulomb
fields. In Section~\ref{sec:many-body}, a uniform one-particle
complementary gap excludes low energies from the omitted continuum
space, and square completion makes coupling to that space enter the
energy quadratically. This gives a lower energy bound on every
antisymmetric form-domain state, uniformly in spin, and a matching
variational upper bound from the mode space. Together they compare the
full continuum infimum with the finite fermion ground energy. Retaining
the direct Coulomb interaction and its counterterm is necessary at the
polynomial scales used here.
Section~\ref{sec:implementation} completes the deterministic computation
of the nuclear positions and thresholds, including all scalar shifts.
